\chapter{Model gromadzenia zapasów}

\section{Wprowadzenie do problemu}

Model gromadzenia zapasów stanowi klasyczne zastosowanie teorii procesów decyzyjnych. W tym rozdziale wykorzystujemy opracowane algorytmy do rozwiązania problemu optymalnego gromadzenia zapasów z dyskontowaniem quasi-hiperbolicznym.

\section{Formalizacja modelu}

Definiujemy model gromadzenia zapasów jako MDP z następującymi składnikami:
\begin{itemize}
\item Stany: poziom zapasów
\item Akcje: wielkość zamówienia  
\item Nagrody: zysk minus koszty przechowywania i zamówienia
\item Dynamika: losowe zużycie zapasów
\end{itemize}

\section{Parametry modelu}

\begin{itemize}
\item Koszt przechowywania jednostki zapasu: $h$
\item Koszt zamówienia: $c$ 
\item Koszt niedoboru: $p$
\item Rozkład popytu: $\{D_t\}$
\end{itemize}

\section{Optymalna polityka}

Pokazujemy, że w modelu z dyskontowaniem quasi-hiperbolicznym optymalna polityka gromadzenia zapasów nie jest spójna czasowo, co ma istotne implikacje praktyczne.

\section{Wyniki numeryczne}

Przedstawiamy wyniki numeryczne dla różnych parametrów modelu, porównując polityki optymalne dla dyskontowania wykładniczego i quasi-hiperbolicznego.

\section{Interpretacja ekonomiczna}

Analizujemy ekonomiczne znaczenie otrzymanych wyników i ich implikacje dla zarządzania zapasami w praktyce.